\documentclass[10pt]{article}
\usepackage[T1]{fontenc}
\usepackage[utf8]{inputenc}
\usepackage[french]{babel}
\usepackage{amsmath,amssymb,mathrsfs,mathpazo}
\usepackage{geometry}
\geometry{a4paper,margin=2cm}
\usepackage{graphicx,array,tabularx,tabularray,colortbl,multirow,multicol,booktabs,enumitem,pifont,ulem,fancybox,diagbox,cancel}
\usepackage{pstricks,pstricks-add,pst-plot,pst-tree,pst-node,pst-text,pst-func,pst-math,pst-bspline,pst-3dplot,pst-eucl,pst-solides3d}
\let\pscalculate\undefined
\usepackage{pst-bezier}
\usepackage{tikz,pgfplots,tkz-tab}
\usetikzlibrary{arrows,arrows.meta,patterns,calc,positioning,trees,decorations.markings,shapes.geometric,shapes.misc,shapes.arrows,intersections,decorations.pathmorphing,decorations.pathreplacing,backgrounds,shadows,babel}
\pgfplotsset{compat=1.18}
\usepgfplotslibrary{fillbetween}
\usepackage[np]{numprint}
\usepackage{listings,makecell,pas-tableur,twoopt}
\let\barre\undefined
\usepackage[pstricks]{tablvar}
\definecolor{darkgreen}{RGB}{0,100,0}
\providecommand{\red}{\color{red}}
\providecommand{\blue}{\color{blue}}
\providecommand{\green}{\color{green}}
\providecommand{\black}{\color{black}}
\providecommand{\white}{\color{white}}
\providecommand{\gray}{\color{gray}}
\providecommand{\cyan}{\color{cyan}}
\providecommand{\magenta}{\color{magenta}}
\providecommand{\yellow}{\color{yellow}}
\providecommand{\violet}{\color{violet}}
\providecommand{\orange}{\color{orange}}
\providecommand{\euro}{\texteuro}
\providecommand{\EUR}{\texteuro}
\providecommand{\ieme}{\textsuperscript{e}}
\providecommand{\degres}{\ensuremath{^\circ}}
\providecommand{\degree}{\ensuremath{^\circ}}
\providecommand{\vv}[1]{\overrightarrow{#1}}
\providecommand{\Coord}[1]{#1}
\providecommand{\decosix}{}
\providecommand{\starredbullet}{\ensuremath{\bullet}}
\pagestyle{empty}
\setlength{\parindent}{0pt}
\usepackage{adjustbox}\DeclareUnicodeCharacter{2236}{:}\DeclareUnicodeCharacter{2190}{\ensuremath{\leftarrow}}
\providecommand{\R}{\mathbb{R}}
\providecommand{\N}{\mathbb{N}}
\providecommand{\Z}{\mathbb{Z}}
\providecommand{\Q}{\mathbb{Q}}
\providecommand{\D}{\mathbb{D}}
\providecommand{\C}{\mathbb{C}}
\providecommand{\e}{{\mathrm{e}}}
\providecommand{\vect}[1]{\overrightarrow{#1}}
\providecommand{\vectt}[1]{\overrightarrow{\mathrm{#1}}}
\providecommand{\vvt}[1]{\overrightarrow{\mathrm{#1}}}
\providecommand{\np}[2][]{{#2~#1}}
\providecommand{\vv}[1]{{\overrightarrow{#1}}}
\renewcommand{\d}{{\mathrm{d}}}
\providecommand{\degres}{{^\circ}}
\providecommand{\degre}{{^\circ}}
\providecommand{\blue}{\color{blue}}
\providecommand{\red}{\color{red}}
\providecommand{\enskip}{\quad}
\providecommand{\llbracket}{\lbrack\!\lbrack}
\providecommand{\rrbracket}{\rbrack\!\rbrack}
\providecommand{\euro}{{\text{€}}}
\providecommand{\no}{{\text{n°}}}
\providecommand{\interval}[2]{{[#1;#2]}}
\providecommand{\intervFO}[2]{{[#1;#2[}}
\providecommand{\dotfill}{\ldots}

\providecommand{\ds}{\displaystyle}
\providecommand{\textding}[1]{#1}
\providecommand{\Oij}{\(\mathrm{O};\vec{\imath},\vec{\jmath}\)}
\providecommand{\Oijk}{\(\mathrm{O};\vec{\imath},\vec{\jmath},\vec{k}\)}
\providecommand{\Ouv}{\(\mathrm{O};\vec{u},\vec{v}\)}
\providecommand{\intervFF}[2]{\ensuremath{[#1;#2]}}
\renewcommand{\euro}{\ifmmode\text{\texteuro}\else\texteuro\fi}
\DeclareUnicodeCharacter{00B0}{\ifmmode^{\circ}\else\textdegree\fi}
\DeclareUnicodeCharacter{20AC}{\ifmmode\text{\texteuro}\else\texteuro\fi}
% Énoncés : Éducation nationale. Transcription et corrigés : APMEP, Denis Vergès et contributeurs. Découpage et classement : CoopMaths, Sébastien Lozano, Éric Elter et contributeurs. Adaptation de la présentation : Maths Explorer.
% CC BY-SA : https://coopmaths.fr/www/about/
% https://coopmaths.fr/alea/static/bac/2024/tex/bac_2024_06_sujet2_metropole_devoile_4_cor.tex
\begin{document}
\noindent\textit{APMEP / CoopMaths. CC BY-SA. Présentation adaptée par Maths Explorer.}\par\medskip
\medskip













\begin{center}
\psset{unit=1cm,arrowsize=2pt 3}
\begin{pspicture*}(-2.5,-1.5)(5.5,3.5)
\psaxes[linewidth=1.25pt,labelFontSize=\scriptstyle]{->}(0,0)(-2.5,-1.5)(5.5,3.5)
\psplot[plotpoints=3000,linewidth=1.25pt,linecolor=red]{-2.5}{5.5}{x 2 add 2.71828 x exp div}
\psplotTangent{0}{4.5}{x 2 add 2.71828 x exp div}
\uput[ul](-2,0){M} \uput[ur](0,2){N} \uput[ur](2,0){P}\uput[d](3,-1){$T$}
\end{pspicture*}

\end{center}

\begin{center}\textbf{Partie A : étude graphique}\end{center}

On répondra aux questions suivantes en utilisant le graphique.
\begin{enumerate}
\item 
	\begin{enumerate}
		\item 
On lit $f(0) = 2$.
		\item Déterminer $f'(0)$.
Le nombre dérivé $f'(0)$ est égal au coefficient directeur de la droite (NP) soit à $\dfrac{0 - 2}{2 - 0} = - 1 = f'(0)$
	\end{enumerate}
\item 
Il semble que $f(-2) = 0$. $S = \{-2\}$.
\item 
Il semble que la fonction est convexe sur l'intervalle $[0~;~+ \infty[$ : sur cet intervalle toutes les tangentes à la courbe représentative $\mathcal{C}_f$ sont sous cette courbe.

Toujours graphiquement sur l'intervalle $[0~;~+ \infty[$ les coefficients directeurs des tangentes à la courbe ($- 1$ et 0 au voisinage de plus l'infini) sont croissants : autrement dit la fonction $f"$ est croissante donc $f''(x) \geqslant 0$
\item 

\begin{center}
\begin{tabularx}{\linewidth}{|*{3}{>{\centering \arraybackslash}X|}}\hline
Courbe 1&Courbe 2&Courbe 3\\ \hline
\psset{unit=0.4cm,arrowsize=2pt 3}
\begin{pspicture*}(-2.5,-8)(8,3)
\psaxes[linewidth=1.25pt,labelFontSize=\scriptstyle]{->}(0,0)(-2.5,-8)(8,3)
\psplot[plotpoints=3000,linewidth=1.25pt,linecolor=red]{-2.5}{5.5}{x 1 add 2.71828 x exp div neg}
\end{pspicture*}&
\psset{unit=0.4cm,arrowsize=2pt 3}
\begin{pspicture*}(-4,-8)(6,3)
\psaxes[linewidth=1.25pt,labelFontSize=\scriptstyle]{->}(0,0)(-4,-8)(6,3)
\psplot[plotpoints=3000,linewidth=1.25pt,linecolor=red]{-4}{5.5}{x 3 add 2.71828 x exp div neg}
\end{pspicture*}
&\psset{unit=0.4cm,arrowsize=2pt 3}
\begin{pspicture*}(-2.5,-8)(8,3)
\psaxes[linewidth=1.25pt,labelFontSize=\scriptstyle]{->}(0,0)(-2.5,-8)(8,3)
\psplot[plotpoints=3000,linewidth=1.25pt,linecolor=red]{-2.5}{5.5}{x 1 sub 2.71828 x exp div neg}
\end{pspicture*}\\ \hline
\end{tabularx}
\end{center}

$\bullet~~$On sait que $f(x) \geqslant 0$ sur $[0~;~+ \infty[$, donc une primitive sur le même intervalle est croissante ce qui élimine la courbe 3 ;

$\bullet~~$Si $F$ est une primitive de $f$ et est représentée par la courbe 1, alors $F'(0) = 0 = f(0) = - 1$ : ceci est faux donc la courbe 1 est éliminée.

Il ne reste que la courbe 2.
\end{enumerate}

\medskip

\begin{center}\textbf{Partie B :  recherche d'une expression algébrique}\end{center}

On admet que la fonction $f$ est de la forme 
\[f(x)= (ax + b)\e^{\lambda x},\]
où $a, b$ et $\lambda$ sont des constantes réelles.



\medskip

\begin{enumerate}
\item 
On a vu que $f(0) = 2 \iff b\e^{\lambda \times 0} = 2 \iff b = 2$.

Donc $f(x) = (ax + 2)\e^{\lambda x}$.
\item 

On a donc $f(x) = (ax + 2)\e^{\lambda x}$.

On sait aussi que $f(-2) = 0 \iff (-2a + 2)\e^{-2\lambda } = 0$ et comme $\e^{-2\lambda }  \ne 0$ on a donc $-2a + 2 = 0 \iff a = 1$.

Donc $f(x) = (x + 2)\e^{\lambda x}$.
\item 
On a vu que $f'(0) = - 1$

La fonction $f$ est dérivable sur $\R$ et sur cet intervalle :

$f'(x) = \e^{\lambda x} + \lambda (x + 2)\e^{\lambda x} = \e^{\lambda x}(1 + \lambda x + 2\lambda )\e^{\lambda x}$.

Donc $f'(0) = - 1 \iff 1  + 2\lambda = - 1 \iff 2\lambda  = - 2 \iff \lambda = - 1 $.

On a donc finalement $f(x) = (x + 2)\e^{- x}$.
\end{enumerate}

\medskip

\begin{center}\textbf{Partie C : étude algébrique}\end{center}

On admet que la fonction $f$ est définie sur $\R$ par 
\[f(x) = (x + 2)\e^{-x}.\]

\smallskip

\begin{enumerate}
\item 
On a $\displaystyle\lim_{x \to -\infty} x + 2 = - \infty$ et $\displaystyle\lim_{x \to -\infty}
\e^{-x} = + \infty$, donc par produit de limites $\displaystyle\lim_{x \to -\infty}f(x) = - \infty$.

\item 
Produit de fonctions dérivables sur $\R$, $f$ est dérivable et sur cet intervalle :

$f'(x) = \e^{-x} - (x + 2)\e^{-x} = \e^{-x}(1 - x - 2) = (- x - 1)\e^{-x}$.

On sait que quel que soit $x \in \R, \: \e^{-x} > 0$ : le signe de $f('x)$ est donc celui de $- x - 1$.

$\bullet -x - 1 > 0 \iff - 1 > x \iff x < - 1$ : sur l'intervalle $]- \infty~;~-1[ \: f'(x) > 0$ : la fonction $f$ est croissante sur cet intervalle ;

$\bullet -x - 1 < 0 \iff - 1 < x \iff x > - 1$ : sur l'intervalle $]-1 ~;~+\infty[\: f'(x) < 0$ : la fonction $f$ est décroissante sur cet intervalle ;

$\bullet -x - 1 = 0 \iff - 1 = x, f'(1) = 0$ ; $f(-1) = (-1 + 2)\e^{1} = \e$ est le maximum de $f$ sur $\R$.
Il reste à calculer la liite en plus l'infini : comme $f(x) = x\e^{-x} + 2\e^{-x}$.

Or on sait que $\displaystyle\lim_{x \to + \infty} x\e^{-x} = 0$ (puissances comparées) et $\displaystyle\lim_{x \to + \infty} 2\e^{-x} = 0$, d'où par somme de limites :
$\displaystyle\lim_{x \to + \infty} f(x) = 0$.

\item 
	\begin{enumerate}
		\item 
$f'$ est elle-même dérivable sur $\R$ et sur cet intervalle :

$f''(x) = -\e^{-x} + (x + 1)\e^{-x} = \e^{-x}(x + 1 - 1) = x \e^{-x}$.

Comme $\e^{-x} > 0$, le signe de $f''(x)$ est celui de $x$, donc :

$\bullet~~$$f$ est convexe sur $[0~;~+ \infty[$ ;

$\bullet~~$$f$ est concave sur $]- \infty~;~0]$ ;
		\item 
-- Donc d'après le résultat précédent la courbe $\mathcal{C}_f$ a un seul point d'inflexion de coordonnées (0~;~2).
	\end{enumerate}
\item 
\[I(t) = \displaystyle\int_{-2}^t (x + 2)\e^{-x}\:\text{d}x.\]

	\begin{enumerate}
		\item 
On pose $u(x) = x + 2$ et $v'(x) = \e^{-x}$, d'où	

$u'(x) = 1$ et $v(x) = - \e^{-x}$.

Toutes ces fonctions étant continues car dérivables on peut intégrer par parties :

$I(t) = \left[-(x + 2)\e^{-x} \right]_{-2}^t + \displaystyle\int_{-2}^t \e^{-x}\:\text{d}x = \left[-(x + 2)\e^{-x}- \e^{-x}\right]_{-2}^t = \left[(- x - 3) \e^{-x} \right]_{-2}^t = (- t - 3) \e^{-t} +1 \e^{2} = \e^{2} - (t + 3)\e^{-t}$.


		\item 
$f$ est positive sur l'intervalle $[-2~;~+ \infty[$ ; on sait qu'alors l'aire de la surface limitée par la courbe $\mathcal{C}_f$, l'axe des abscisses et les droites d'équations $x = - 2$ et $x = t$ est égale à l'intégrale $I(t)$.

Or quand $t \to + \infty$, la surface n'est pas limitée à droite alors que l'intégrale l'est elle puisqu'on a $\displaystyle\lim_{t \to + \infty} (t + 3)\e^{- t} = 0$, donc $\displaystyle\lim_{t \to + \infty} I(t) = \e^2$.
	\end{enumerate}
\end{enumerate}

\end{document}
